% संपूर्ण मराठी ग्रंथातील प्रकरण 48: मोडल टॅब्लो
% हा संपादनयोग्य स्रोतखंड आहे. संकलनासाठी संपूर्ण स्रोतसंग्रहातील
% अक्षररूपे, पूर्वभाग, चिन्हव्याख्या आणि आकृती-स्रोत आवश्यक आहेत.
% मूळ: Open Logic Project; मजकूर आणि रूपांतर CC BY 4.0.
% यंत्रानुवाद, दुरुस्ती आणि तपासणी: OpenAI Codex —
% GPT-5.6 Sol आणि GPT-6 Sol, दोन्ही Ultra effort.
% दुरुस्ती-पुनरावलोकन: GPT-6.1 Sol, Ultra effort.
\chapter{मोडल टॅब्लो}\label{nml:tab:chap}









\section{प्रस्तावना}\label{nml:tab:int:sec}

टॅब्लो म्हणजे चिन्हांकित सूत्रांचे (खाली शाखा फुटणारे) काही वृक्ष.
चिन्हांकित सूत्र म्हणजे सत्यतामूल्याचे चिन्ह
($\True$ किंवा $\False$) आणि वाक्य
यांची जोडी:
\[\sFmla{\True}{\varphi} \text{किंवा} \sFmla{\False}{\varphi}.
\]
टॅब्लोची सुरुवात काही \emph{गृहीतकां}पासून
होते. पुढील प्रत्येक चिन्हांकित सूत्र हे अनुमानाचा
एखादा नियम लावून मिळते. काही नियम वृक्षाच्या
टोकावर एक किंवा अधिक चिन्हांकित सूत्रे जोडतात;
इतर नियम दोन नवी टोके जोडून दोन शाखा तयार करतात.
नियम लावल्यावर मिळणाऱ्या चिन्हांकित सूत्रे मधील
सूत्र, नियम ज्या चिन्हांकित सूत्राला लावला
त्यातील सूत्रापेक्षा कमी गुंतागुंतीचे असते.
एखाद्या शाखेत $\sFmla{\True}{\varphi}$ आणि
$\sFmla{\False}{\varphi}$ ही दोन्ही असतील, तर ती शाखा
\emph{बंद} आहे असे म्हणतो. टॅब्लो मधील प्रत्येक
शाखा बंद असेल, तर संपूर्ण टॅब्लो बंद असतो.
बंद टॅब्लो ही अशी निष्पत्ती असते की
तिच्यावरून त्या टॅब्लोची सुरुवात ज्या
चिन्हांकित सूत्रांनी केली तो संच अपूर्ततायोग्य
आहे हे दिसते. यावरून $\Proves$ संबंधाची व्याख्या
करता येते: $\Gamma \Proves \varphi$ तेव्हा आणि तेव्हाच
$\Gamma_0 = \{\psi_1,
\dots, \psi_n\} \subseteq \Gamma$ असा एखादा सांत
संच असतो की पुढील गृहीतकांसाठी बंद टॅब्लो
मिळतो:
\[\{\sFmla{\False}{\varphi}, \sFmla{\True}{\psi_1}, \dots, \sFmla{\True}{\psi_n}\}.
\]

मोडल तर्कशास्त्रांसाठी चिन्हांकित सूत्र ही संकल्पना विस्तारित करावी लागते,
तसेच~$\Box$ आणि
    $\Diamond$ यांसाठी नियम जोडावे लागतात.
चिन्हाबरोबर ($\True$ किंवा $\False$) मोडल
टॅब्लो मधील सूत्रे ना
\emph{पूर्वचिन्हे}~$\sigma$ ही असतात.
ही पूर्वचिन्हे धन पूर्णांकांच्या रिकाम्या नसलेल्या
क्रमिका असतात; म्हणजे $\sigma \in (\PosInt)^* \setminus
\{\emptyseq\}$. अशी पूर्वचिन्हे लिहिताना भोवतालचे
$\tuple{\ }$ वगळतो आणि वेगवेगळ्या घटक
मध्ये $,$ ऐवजी $.$ वापरतो.
$\sigma$ हे पूर्वचिन्ह असेल, तर $\sigma.n$ म्हणजे
$\sigma \concat \tuple{n}$; उदाहरणार्थ
$\sigma = 1.2.1$ असेल, तर $\sigma.3$ म्हणजे
$1.2.1.3$. उदाहरणार्थ,
\[\sFmla{\True}{\Box \varphi \lif \varphi}[1.2]
\]
हे \emph{पूर्वचिन्हयुक्त चिन्हांकित सूत्र}
(किंवा संक्षेपाने \emph{पूर्वचिन्हयुक्त सूत्र}) आहे.

सहज अर्थाने, टॅब्लोच्या एखाद्या शाखेवरील
सूत्रे पूर्ण करू शकेल अशा प्रतिमानातील
जगाचे नाव पूर्वचिन्ह देते. आणि $\sigma$ हे
एखाद्या जगाचे नाव असेल, तर $\sigma.n$ हे
म्हणजे~$\sigma$ ने दर्शवलेल्या जगापासून प्राप्य
असलेल्या जगाचे नाव देते.













\section{\Ax{K} साठीचे नियम}\label{nml:tab:rul:sec}

नेहमीच्या विधानीय कारकांचे नियम नेहमीच्या विधानीय
चिन्हांकित टॅब्लो मधील नियमांसारखेच आहेत;
फक्त त्यांना पूर्वचिन्हे जोडली जातात. प्रत्येक वेळी,
$\sFmla{S}{\varphi}[\sigma]$ या चिन्हांकित सूत्र
ला नियम लावल्याने मिळणाऱ्या नव्या सूत्रे
वरही~$\sigma$ हेच पूर्वचिन्ह असते. हे सहज लक्षात
येते: उदाहरणार्थ, $\varphi \land \psi$ हे~$\sigma$ ने
नाव दिलेल्या जगात सत्य असेल, तर $\varphi$ आणि $\psi$
ही दोन्ही~$\sigma$ येथे सत्य असतात (इतर कोणत्याही
जगात असतीलच असे नाही). विधानीय नियम
\ref{nml:tab:rul:tab:prop-rules} मध्ये दिले आहेत.

\begin{table}
  \[\def\arraystretch{3}\begin{array}{|c|c|}
    \hline
    \AxiomC{\sFmla{\True}{\lnot \varphi}[\sigma]}
    \RightLabel{\TRule{\True}{\lnot}}
    \UnaryInfC{\sFmla{\False}{\varphi}[\sigma]}
    \DisplayProof
    &
    \AxiomC{\sFmla{\False}{\lnot \varphi}[\sigma]}
    \RightLabel{\TRule{\False}{\lnot}}
    \UnaryInfC{\sFmla{\True}{\varphi}[\sigma]}
    \DisplayProof
    \\[1ex]
    \hline
    \AxiomC{\sFmla{\True}{\varphi \land \psi}[\sigma]}
    \RightLabel{\TRule{\True}{\land}}
    \UnaryInfC{\sFmla{\True}{\varphi}[\sigma]}
    \noLine
    \UnaryInfC{\sFmla{\True}{\psi}[\sigma]}
    \DisplayProof
    &
    \AxiomC{\sFmla{\False}{\varphi \land \psi}[\sigma]}
    \RightLabel{\TRule{\False}{\land}}
    \UnaryInfC{$\sFmla{\False}{\varphi}[\sigma] \quad \mid \quad
      \sFmla{\False}{\psi}[\sigma]$}
    \DisplayProof
    \\[2ex]
    \hline
    \AxiomC{\sFmla{\True}{\varphi \lor \psi}[\sigma]}
    \RightLabel{\TRule{\True}{\lor}}
    \UnaryInfC{$\sFmla{\True}{\varphi}[\sigma] \quad \mid \quad
      \sFmla{\True}{\psi}[\sigma]$}
    \DisplayProof
    &
    \AxiomC{\sFmla{\False}{\varphi \lor \psi}[\sigma]}
    \RightLabel{\TRule{\False}{\lor}}
    \UnaryInfC{\sFmla{\False}{\varphi}[\sigma]}
    \noLine
    \UnaryInfC{\sFmla{\False}{\psi}[\sigma]}
    \DisplayProof
    \\[2ex]
    \hline
    \AxiomC{\sFmla{\True}{\varphi \lif \psi}[\sigma]}
    \RightLabel{\TRule{\True}{\lif}}
    \UnaryInfC{$\sFmla{\False}{\varphi}[\sigma] \quad \mid
      \quad \sFmla{\True}{\psi}[\sigma]$}
    \DisplayProof
    &
    \AxiomC{\sFmla{\False}{\varphi \lif \psi}[\sigma]}
    \RightLabel{\TRule{\False}{\lif}}
    \UnaryInfC{\sFmla{\True}{\varphi}[\sigma]}
    \noLine
    \UnaryInfC{\sFmla{\False}{\psi}[\sigma]}
    \DisplayProof
    \\[2ex]
    \hline
  \end{array}\]
  \caption{विधानीय कारकांसाठी पूर्वचिन्हयुक्त
    टॅब्लो नियम}
  \label{nml:tab:rul:tab:prop-rules}
\end{table}

शाखा बंद होण्याची अट नेहमीच्या टॅब्लो
प्रमाणेच आहे; पण येथे सूत्रे बरोबर त्यांची
पूर्वचिन्हेही जुळली पाहिजेत. म्हणजे एखाद्या शाखेत
पुढील दोन्ही असतील, तर ती बंद असते:
\[\sFmla{\True}{\varphi}[\sigma] \quad\text{आणि}\quad \sFmla{\False}{\varphi}[\sigma]
\]
येथे $\sigma$ हे एखादे पूर्वचिन्ह आणि $\varphi$ हे
सूत्र आहे.

गृहीतके मांडण्याचे नियमही नेहमीच्या टॅब्लो
प्रमाणेच आहेत; मात्र गृहीतकांसाठी आपण नेहमी
$1$ हे पूर्वचिन्ह वापरतो. (सर्व गृहीतकांसाठी एकच
पूर्वचिन्ह वापरले, तर नेमके कोणते पूर्वचिन्ह
वापरले याने फरक पडत नाही.) म्हणून, उदाहरणार्थ,
\[\psi_1, \dots, \psi_n \Proves \varphi
\]
असे तेव्हा आणि तेव्हाच म्हणतो जेव्हा पुढील
गृहीतकांसाठी बंद टॅब्लो असतो:
\[\sFmla{\True}{\psi_1}[1], \dots, \sFmla{\True}{\psi_n}[1],
\sFmla{\False}{\varphi}[1].
\]

मोडल कारक~$\Box$
    आणि~$\Diamond$
यांसाठी,~$\sigma$ हे पूर्वचिन्ह असलेल्या सूत्र
ला नियम लावल्यावर निष्कर्षाचे पूर्वचिन्ह
$\sigma.n$ असते. पण कोणता $n$ चालेल हे
चिन्ह~$\True$ आहे की~$\False$ यावर अवलंबून असते.

$\TRule{\True}{\Box}$ हा नियम
  $\sFmla{\True}{\Box \varphi}[\sigma]$ असलेल्या शाखेत
  $\sFmla{\True}{\varphi}[\sigma.n]$ जोडतो.
    त्याचप्रमाणे, $\TRule{\False}{\Diamond}$
  हा नियम $\sFmla{\False}{\Diamond \varphi}[\sigma]$
  असलेल्या शाखेत $\sFmla{\False}{\varphi}[\sigma.n]$
  जोडतो.
हे नियम
लावताना $\sigma.n$ हे पूर्वचिन्ह संबंधित शाखेत
\emph{आधीपासूनच} आलेले असले पाहिजे. अशा
पूर्वचिन्हाला त्या शाखेवर ‘वापरलेले’ म्हणू.

$\TRule{\False}{\Box}$ हा नियम
  $\sFmla{\False}{\Box \varphi}[\sigma]$ असलेल्या शाखेत
  $\sFmla{\False}{\varphi}[\sigma.n]$ जोडतो.
    त्याचप्रमाणे, $\TRule{\True}{\Diamond}$
  हा नियम $\sFmla{\True}{\Diamond \varphi}[\sigma]$
  असलेल्या शाखेत $\sFmla{\True}{\varphi}[\sigma.n]$
  जोडतो.
परंतु हे नियम
लावताना $\sigma.n$ हे पूर्वचिन्ह संबंधित शाखेत
\emph{आधीपासून आलेले नसले} पाहिजे. अशा
पूर्वचिन्हांना त्या शाखेसाठी ‘नवी’ म्हणतो.

हे नियम \ref{nml:tab:rul:tab:rules-K} मध्ये दिले आहेत.

\begin{table}
  \begin{center}
    \def\arraystretch{3}\def\fCenter{}
    \begin{tabular}{|c|c|}
    \hline

      \AxiomC{\sFmla{\True}{\Box \varphi}[\sigma]}
      \RightLabel{\TRule{\True}{\Box}}
      \UnaryInfC{\sFmla{\True}{\varphi}[\sigma.n]}
      \DisplayProof
      &
      \AxiomC{\sFmla{\False}{\Box \varphi}[\sigma]}
      \RightLabel{\TRule{\False}{\Box}}
      \UnaryInfC{\sFmla{\False}{\varphi}[\sigma.n]}
      \DisplayProof\\
      $\sigma.n$ वापरलेले आहे & $\sigma.n$ नवे आहे
      \\[1ex]
      \hline
      \AxiomC{\sFmla{\True}{\Diamond \varphi}[\sigma]}
      \RightLabel{\TRule{\True}{\Diamond}}
      \UnaryInfC{\sFmla{\True}{\varphi}[\sigma.n]}
      \DisplayProof
      &
      \AxiomC{\sFmla{\False}{\Diamond \varphi}[\sigma]}
      \RightLabel{\TRule{\False}{\Diamond}}
      \UnaryInfC{\sFmla{\False}{\varphi}[\sigma.n]}
      \DisplayProof\\
      $\sigma.n$ नवे आहे & $\sigma.n$ वापरलेले आहे
      \\[1ex]
      \hline
    \end{tabular}
  \end{center}
  \caption{\Ax{K} साठीचे मोडल नियम.}
  \label{nml:tab:rul:tab:rules-K}
\end{table}
\TRule{\True}{\Box}
या नियमासाठी पूर्वचिन्ह वापरलेले असणे आवश्यक आहे.
ही अट नसेल, तर पुढील टॅब्लो बंद असल्याचे
आपण चुकीने मानू:
  \begin{oltableau}
    [\pFmla{\True}{\Box \formula{A}}{1}, just = \TAss
      [\pFmla{\False}{\Diamond \formula{A}}{1}, just = \TAss
        [\pFmla{\True}{\formula{A}}{1.1}, just = {\TRule{\True}{\Box}[1]}
          [\pFmla{\False}{\formula{A}}{1.1},
            just = {\TRule{\False}{\Diamond}[2]}, close]
        ]
      ]
    ]
\end{oltableau}
पण $\Box \formula{A} \Entails/ \Diamond \formula{A}$;
म्हणून असे झाल्यास आपली प्रमाणपद्धती अप्रमाण असेल.
त्याचप्रमाणे, $\Diamond \formula{A} \Entails/
\Box \formula{A}$; परंतु
\TRule{\False}{\Box}
या नियमासाठी पूर्वचिन्ह नवे असण्याची अट नसेल,
तर पुढील टॅब्लो बंद होईल:
  \begin{oltableau}
    [\pFmla{\True}{\Diamond \formula{A}}{1}, just = \TAss,name=one
      [\pFmla{\False}{\Box \formula{A}}{1}, just = \TAss, name=two
        [\pFmla{\True}{\formula{A}}{1.1}, just = {\TRule{\True}{\Diamond}[1]}
          [\pFmla{\False}{\formula{A}}{1.1},
            just = {\TRule{\False}{\Box}[2]}, close]
        ]
      ]
    ]
  \end{oltableau}
\section{\Log{K} साठी टॅब्लो}\label{nml:tab:prk:sec}
\begin{ex}
  $\Proves (\Box\varphi \land \Box\psi)
  \lif \Box (\varphi \land \psi)$ हे दाखवणारा बंद टॅब्लो
  पुढीलप्रमाणे आहे.
  \begin{oltableau}
    [\pFmla{\False}{(\Box\formula{A} \land \Box\formula{B}) \lif
        \Box (\formula{A} \land \formula{B})}{1},
      just =\TAss
      [\pFmla{\True}{\Box\formula{A} \land \Box\formula{B}}{1},
        just = {\TRule{\False}{\lif}[1]}
        [\pFmla{\False}{\Box(\formula{A} \land \formula{B})}{1},
          just = {\TRule{\False}{\lif}[1]}
          [\pFmla{\True}{\Box\formula{A}}{1},
            just = {\TRule{\True}{\land}[2]}
            [\pFmla{\True}{\Box\formula{B}}{1},
              just = {\TRule{\True}{\land}[2]}
              [\pFmla{\False}{\formula{A} \land \formula{B}}{1.1},
                just = {\TRule{\False}{\Box}[3]}
                [\pFmla{\False}{\formula{A}}{1.1},
                  just = {\TRule{\False}{\land}[6]}
                  [\pFmla{\True}{\formula{A}}{1.1},
                    just= {\TRule{\True}{\Box}[4]}, close]]
                [\pFmla{\False}{\formula{B}}{1.1},
                  just = {\TRule{\False}{\land}[6]}
                  [\pFmla{\True}{\formula{B}}{1.1},
                    just= {\TRule{\True}{\Box}[5]}, close]]
              ]
            ]
          ]
        ]
      ]
    ]
  \end{oltableau}
\end{ex}
\begin{ex}
  $\Proves \Diamond(\varphi \lor \psi)
  \lif (\Diamond \varphi \lor \Diamond \psi)$ हे दाखवणारा
  बंद टॅब्लो पुढीलप्रमाणे आहे:
  \begin{oltableau}
    [\pFmla{\False}{\Diamond(\formula{A} \lor \formula{B}) \lif
        (\Diamond \formula{A} \lor \Diamond \formula{B})}{1},
      just =\TAss
      [\pFmla{\True}{\Diamond(\formula{A} \lor \formula{B})}{1},
        just = {\TRule{\False}{\lif}[1]}
        [\pFmla{\False}{\Diamond\formula{A} \lor \Diamond\formula{B}}{1},
          just = {\TRule{\False}{\lif}[1]}
          [\pFmla{\False}{\Diamond\formula{A}}{1},
            just = {\TRule{\False}{\lor}[3]}
            [\pFmla{\False}{\Diamond\formula{B}}{1},
              just = {\TRule{\False}{\lor}[3]}
              [\pFmla{\True}{\formula{A} \lor \formula{B}}{1.1},
                just = {\TRule{\True}{\Diamond}[2]},
                [\pFmla{\True}{\formula{A}}{1.1},
                  just = {\TRule{\True}{\lor}[6]}
                  [\pFmla{\False}{\formula{A}}{1.1},
                    just= {\TRule{\False}{\Diamond}[4]}, close]]
                [\pFmla{\True}{\formula{B}}{1.1},
                  just = {\TRule{\True}{\lor}[6]}
                  [\pFmla{\False}{\formula{B}}{1.1},
                    just= {\TRule{\False}{\Diamond}[5]}, close]]
              ]
            ]
          ]
        ]
      ]
    ]
  \end{oltableau}
\end{ex}
\begin{prob}
  पुढील सूत्रे साठी~$\Log{K}$ मध्ये बंद
  टॅब्लो शोधा:
  \begin{enumerate}
    \item $\Box \lnot p \lif \Box(p \lif q)$
    \item $(\Box p \lor \Box q) \lif \Box(p \lor q)$
    \item $\Diamond p \lif \Diamond(p \lor q)$
    \item $\Box(p \land q) \lif \Box p$
  \end{enumerate}
\end{prob}
\section{\Log{K} साठी निर्दोषता}\label{nml:tab:sou:sec}
\begin{editorial}
  निर्दोषतेचा हा पुरावा अभिजात विधानीय तर्कशास्त्राचा
  निर्दोषतेचा पुरावा पुन्हा मांडतो; म्हणजे सर्व काही
  सुरुवातीपासून सिद्ध करतो. स्वयंपूर्ण पुरावा हवा
  असेल, तर हे योग्य आहे. नेहमीच्या टॅब्लोची निर्दोषता
  तुम्ही आधी पाहिली असेल, तर येथे पुनरुक्ती वाटेल.
  पुढे स्वयंपूर्ण आवृत्ती आणि अमोडल प्रकरणावर आधारलेली
  आवृत्ती यांपैकी एक निवडता यावी, अशी योजना आहे.
\end{editorial}
\begin{explain}
  पूर्वचिन्हयुक्त टॅब्लो निर्दोष आहेत हे
  दाखवण्यासाठी, पुढील गृहीतकांचा
  \[
  \sFmla{\True}{\psi_1}[1], \dots, \sFmla{\True}{\psi_n}[1], \sFmla{\False}{\varphi}[1]
  \]
  बंद टॅब्लो असेल, तर $\psi_1, \dots, \psi_n
  \Entails \varphi$ हे सिद्ध करावे लागेल. व्यत्यास सिद्ध
  करणे सोपे आहे: एखाद्या $\mModel{M}$ मध्ये
  $w$ या जगात सर्व $i=1$, \dots,~$n$ साठी
  $\mSat{M}{\psi_i}[w]$, पण $\mSat/{M}{\varphi}[w]$
  असेल, तर कोणताही टॅब्लो बंद होऊ शकत नाही.
  अशा प्रतिदृष्टांत प्रतिमानावरून टॅब्लोची सुरुवातीची
  गृहीतके पूर्ततायोग्य आहेत हे दिसते. पुराव्याची युक्ती
  अशी: टॅब्लोच्या शाखेवरील सर्व
  पूर्वचिन्हयुक्त सूत्रे पूर्ततायोग्य असतील, तर
  कोणताही नियम लावल्यावर तयार होणाऱ्या विस्तारित
  शाखांपैकी किमान एक शाखा पूर्ततायोग्य असते. बंद शाखा
  अपूर्ततायोग्य असतात. म्हणून पूर्वचिन्हयुक्त सूत्रे
  च्या पूर्ततायोग्य संचासाठीच्या कोणत्याही टॅब्लो
  मध्ये किमान एक खुली शाखा असली पाहिजे.

  मोडल प्रकरणात ही युक्ती वापरण्यासाठी ‘पूर्ततायोग्य’
  ही व्याख्या मोडल प्रतिमाने आणि पूर्वचिन्हे यांच्यापर्यंत
  विस्तारित करावी लागते. त्यानंतर पुरावा सरळ आहे.
\end{explain}

\begin{defn}
  $P$ हा पूर्वचिन्हांचा एखादा संच असू द्या; म्हणजे
  $P \subseteq (\PosInt)^* \setminus
  \{\emptyseq\}$. तसेच $\mModel{M}$ हे प्रतिमान असू द्या.
  $\sigma$ आणि $\sigma.n$ ही दोन्ही~$P$ मध्ये असताना
  नेहमी $Rf(\sigma)f(\sigma.n)$ असेल, तर
  $f\colon P \to W$ या फलनाला $\mModel{M}$
  मधील~$P$ चे \emph{अर्थनिर्धारण} म्हणतात.

  पूर्वचिन्हांच्या~$P$ या संचाच्या अर्थनिर्धारणाच्या
  संदर्भात पुढील व्याख्या करता येतात:
  \begin{enumerate}
  \item $\mModel{M}$ हे $\sFmla{\True}{\varphi}[\sigma]$
    पूर्ण करते तेव्हा आणि तेव्हाच
    $\mSat{M}{\varphi}[f(\sigma)]$.
  \item $\mModel{M}$ हे $\sFmla{\False}{\varphi}[\sigma]$
    पूर्ण करते तेव्हा आणि तेव्हाच
    $\mSat/{M}{\varphi}[f(\sigma)]$.
  \end{enumerate}
\end{defn}

\begin{defn}
  $\Gamma$ हा पूर्वचिन्हयुक्त सूत्रांचा
  संच असू द्या आणि त्यात येणाऱ्या पूर्वचिन्हांचा
  संच $P(\Gamma)$ असू द्या. $f$ हे $\mModel{M}$
  मधील~$P(\Gamma)$ चे अर्थनिर्धारण असेल, तर
  $\mModel{M}$ ने~$f$ च्या संदर्भात~$\Gamma$
  पूर्ण करणे, म्हणजे $\mSat{M}{\Gamma}[f]$,
  याचा अर्थ $\mModel{M}$ ने~$f$ च्या संदर्भात
  $\Gamma$ मधील प्रत्येक पूर्वचिन्हयुक्त सूत्र
  पूर्ण करणे असा आहे. $\mSat{M}{\Gamma}[f]$
  होईल असे एखादे प्रतिमान~$\mModel{M}$ आणि
  $P(\Gamma)$ चे अर्थनिर्धारण~$f$ अस्तित्वात असेल,
  तर आणि तरच $\Gamma$ हा संच \emph{पूर्ततायोग्य} असतो.
\end{defn}

\begin{prop}
  एखाद्या सूत्र~$\varphi$ आणि पूर्वचिन्ह~$\sigma$
  यांसाठी $\Gamma$ मध्ये
  $\sFmla{\True}{\varphi}[\sigma]$ आणि
  $\sFmla{\False}{\varphi}[\sigma]$ ही दोन्ही असतील,
  तर $\Gamma$ अपूर्ततायोग्य आहे.
\end{prop}

\begin{proof}
  $\mSat{M}{\varphi}[f(\sigma)]$ आणि
  $\mSat/{M}{\varphi}[f(\sigma)]$ ही दोन्ही सत्य होतील
  असे प्रतिमान~$\mModel{M}$ आणि $P(\Gamma)$ चे
  अर्थनिर्धारण~$f$ असू शकत नाही.
\end{proof}

\begin{thm}[निर्दोषता]
  \label{nml:tab:sou:thm:tableau-soundness}
  $\Gamma$ साठी बंद टॅब्लो असेल, तर
  $\Gamma$ अपूर्ततायोग्य आहे.
\end{thm}

\begin{proof}
टॅब्लोवरील एखादी शाखा पूर्ततायोग्य असणे म्हणजे
तिच्यावरील चिन्हांकित सूत्रांचा संच पूर्ततायोग्य असणे.
आणि टॅब्लोमध्ये किमान एक पूर्ततायोग्य शाखा
असेल, तर त्याला पूर्ततायोग्य म्हणू.
येथे नियमानुसार तयार केलेल्या टॅब्लोंचाच विचार करतो. सुरुवातीच्या
गृहीतकांचे पूर्वचिन्ह $1$ असते आणि प्रत्येक नवे पूर्वचिन्ह आधी
वापरलेल्या पूर्वचिन्हाच्या शेवटी एक संख्या जोडून तयार होते.
म्हणून शाखेवरील पूर्वचिन्हांचा संच प्रत्येक रिक्त नसलेल्या आरंभीच्या
खंडाखाली बंद असतो.

आपण पुढील विधान दाखवू: पूर्ततायोग्य टॅब्लो ला
अनुमानाचा कोणताही नियम लावून विस्तारित केले, तरी
मिळणारा टॅब्लो पूर्ततायोग्य असतो. यावरून प्रमेय
सिद्ध होईल. कारण बंद टॅब्लो हा
सुरुवातीला फक्त~$\Gamma$ मधील गृहीतके असलेल्या
टॅब्लो ला अनुमानाचे नियम लावत तयार होतो.
म्हणून $\Gamma$ पूर्ततायोग्य असेल, तर त्यासाठीचा
प्रत्येक टॅब्लो पूर्ततायोग्य राहील. पण बंद
टॅब्लो पूर्ततायोग्य असू शकत नाही: त्याच्या सर्व
शाखा बंद असतात, आणि बंद शाखा अपूर्ततायोग्य असतात.

आपल्याकडे पूर्ततायोग्य टॅब्लो आहे असे समजू;
म्हणजे किमान एक पूर्ततायोग्य शाखा असलेला टॅब्लो
आहे. अनुमानाचा
नियम लावल्याने शाखेवर चिन्हांकित सूत्रे
जोडली जातात किंवा शाखेच्या दोन शाखा फुटतात.
नियमामुळे न बदलणारी एखादी पूर्ततायोग्य शाखा
टॅब्लोमध्ये असेल, तर विस्तारित
टॅब्लोमध्येही ती पूर्ततायोग्य राहते. म्हणून
नियम पूर्ततायोग्य शाखेला लावला जातो त्या
प्रकरणाचा विचार पुरेसा आहे.

त्या शाखेवरील चिन्हांकित सूत्रांचा संच
$\Gamma$ असू द्या आणि नियम
$\sFmla{S}{\varphi}[\sigma] \in \Gamma$ या
चिन्हांकित सूत्राला लावला जातो असे समजू.
नियमाने शाखा फुटत नसेल, तर $\Gamma$ मध्ये
नियमाचे निष्कर्ष जोडून मिळालेली विस्तारित
शाखा अजूनही पूर्ततायोग्य आहे हे दाखवावे लागते.
शाखा फुटत असेल, तर मिळणाऱ्या दोन शाखांपैकी
किमान एक पूर्ततायोग्य आहे हे दाखवावे लागते.

प्रथम शाखा न फुटणाऱ्या नियमांचा विचार करू.
\begin{enumerate}
\item $\sFmla{\True}{\lnot \psi}[\sigma] \in \Gamma$
  याला
  $\TRule{\True}{\lnot}$ लावून शाखा विस्तारित
  केली. आता शाखेवरील चिन्हांकित सूत्रांचा
  संच $\Gamma \cup
  \{\sFmla{\False}{\psi}[\sigma]\}$ आहे.
  $\mSat{M}{\Gamma}[f]$ असे समजू. विशेषतः
  $\mSat{M}{\lnot \psi}[f(\sigma)]$; म्हणून
  $\mSat/{M}{\psi}[f(\sigma)]$. म्हणजे
  $\mModel{M}$ हे~$f$ च्या संदर्भात
  $\sFmla{\False}{\psi}[\sigma]$ पूर्ण करते.
\item $\sFmla{\False}{\lnot \psi}[\sigma] \in \Gamma$
  याला
  $\TRule{\False}{\lnot}$ लावून शाखा विस्तारित
  केली: सरावासाठी.
\item $\sFmla{\True}{\psi \land \chi}[\sigma] \in \Gamma$
  याला
  $\TRule{\True}{\land}$ लावून शाखा विस्तारित
  केली. यात शाखेवर
  $\sFmla{\True}{\psi}[\sigma]$ आणि
  $\sFmla{\True}{\chi}[\sigma]$ ही दोन नवी
  चिन्हांकित सूत्रे मिळतात.
  $\mSat{M}{\Gamma}[f]$ असे समजू; विशेषतः
  $\mSat{M}{\psi \land \chi}[f(\sigma)]$.
  म्हणून $\mSat{M}{\psi}[f(\sigma)]$ आणि
  $\mSat{M}{\chi}[f(\sigma)]$. याचा अर्थ
  $\mModel{M}$ हे~$f$ च्या संदर्भात
  $\sFmla{\True}{\psi}[\sigma]$ आणि
  $\sFmla{\True}{\chi}[\sigma]$ ही दोन्ही
  पूर्ण करते.
\item $\sFmla{\False}{\psi \lor \chi}[\sigma] \in \Gamma$
  याला
  $\TRule{\False}{\lor}$ लावून शाखा विस्तारित
  केली: सरावासाठी.
\item $\sFmla{\False}{\psi \lif \chi}[\sigma] \in \Gamma$
  याला
  $\TRule{\False}{\lif}$ लावून शाखा विस्तारित
  केली. यात शाखेवर
  $\sFmla{\True}{\psi}[\sigma]$ आणि
  $\sFmla{\False}{\chi}[\sigma]$ ही दोन नवी
  चिन्हांकित सूत्रे मिळतात.
  $\mSat{M}{\Gamma}[f]$ असे समजू; विशेषतः
  $\mSat/{M}{\psi \lif \chi}[f(\sigma)]$.
  म्हणून $\mSat{M}{\psi}[f(\sigma)]$ आणि
  $\mSat/{M}{\chi}[f(\sigma)]$. याचा अर्थ
  $\mModel{M}, f$ हे
  $\sFmla{\True}{\psi}[\sigma]$ आणि
  $\sFmla{\False}{\chi}[\sigma]$ ही दोन्ही
  पूर्ण करते.


\item $\sFmla{\True}{\Box \psi}[\sigma] \in \Gamma$
  याला
  $\TRule{\True}{\Box}$ लावून शाखा विस्तारित केली:
  शाखेवर
    $\sFmla{\True}{\psi}[\sigma.n]$ हे नवे
    चिन्हांकित सूत्र मिळते. येथे
    $\sigma.n \in P(\Gamma)$, कारण
    $\sigma.n$ वापरलेले असले पाहिजे.
    $\mSat{M}{\Gamma}[f]$ असे समजू; विशेषतः
    $\mSat{M}{\Box \psi}[f(\sigma)]$.
    $f$ हे पूर्वचिन्हांचे अर्थनिर्धारण आहे आणि
    $\sigma$, $\sigma.n \in P(\Gamma)$,
    म्हणून $Rf(\sigma)f(\sigma.n)$.
    त्यामुळे $\mSat{M}{\psi}[f(\sigma.n)]$;
    म्हणजे $\mModel{M}, f$ हे
    $\sFmla{\True}{\psi}[\sigma.n]$ पूर्ण करते.

\item $\sFmla{\False}{\Box \psi}[\sigma] \in \Gamma$
  याला
  $\TRule{\False}{\Box}$ लावून शाखा विस्तारित केली:
  शाखेवर
    $\sFmla{\False}{\psi}[\sigma.n]$ हे नवे
    चिन्हांकित सूत्र मिळते. येथे
    $\sigma.n$ हे शाखेवरील नवे पूर्वचिन्ह,
    म्हणजे $\sigma.n \notin P(\Gamma)$.
    $\Gamma$ पूर्ततायोग्य असल्याने,
    $\mSat{M}{\Gamma}[f]$ होईल असे
    $\mModel{M}$ आणि $P(\Gamma)$ चे
    अर्थनिर्धारण~$f$ असते; विशेषतः
    $\mSat/{M}{\Box \psi}[f(\sigma)]$.
    $\Gamma \cup
    \{\sFmla{\False}{\psi}[\sigma.n]\}$
    पूर्ततायोग्य आहे हे दाखवायचे आहे.
    यासाठी $P(\Gamma) \cup \{\sigma.n\}$
    चे अर्थनिर्धारण पुढीलप्रमाणे ठरवू:

  $\mSat/{M}{\Box \psi}[f(\sigma)]$ असल्याने,
  $Rf(\sigma)w$ आणि $\mSat/{M}{\psi}[w]$
  होईल असे $w \in W$ असते. $f'$ हे~$f$
  प्रमाणेच असू द्या, फक्त
  $f'(\sigma.n) = w$ असे ठरवा.
  $f'(\sigma) = f(\sigma)$ आणि
  $Rf(\sigma)w$ असल्याने
  $Rf'(\sigma)f'(\sigma.n)$.
  पूर्वचिन्हांच्या आरंभीच्या खंडांखालील बंदतेमुळे नव्या
  $\sigma.n$ पासून वाढवलेले कोणतेही पूर्वचिन्ह आधी वापरलेले
  नसते. त्यामुळे अर्थनिर्धारणासाठी तपासायची नवी जोडी फक्त
  $(\sigma,\sigma.n)$ हीच आहे; म्हणून $f'$
  हे $P(\Gamma) \cup \{\sigma.n\}$ चे
  अर्थनिर्धारण आहे. स्पष्टपणे
  $\mSat/{M}{\psi}[f'(\sigma.n)]$.
  प्रत्येक $\sigma' \in P(\Gamma)$ साठी
  $f(\sigma') = f'(\sigma')$ असल्याने
  $\mSat{M}{\Gamma}[f']$. म्हणून
  $\mModel{M}, f'$ हे $\Gamma \cup
  \{\sFmla{\False}{\psi}[\sigma.n]\}$ पूर्ण करते.



\item $\sFmla{\True}{\Diamond \psi}[\sigma] \in \Gamma$
  याला
  $\TRule{\True}{\Diamond}$ लावून शाखा विस्तारित
  केली: शाखेवर
    $\sFmla{\True}{\psi}[\sigma.n]$ हे नवे
    चिन्हांकित सूत्र मिळते. येथे
    $\sigma.n$ हे शाखेवरील नवे पूर्वचिन्ह,
    म्हणजे $\sigma.n \notin P(\Gamma)$.
    $\Gamma$ पूर्ततायोग्य असल्याने,
    $\mSat{M}{\Gamma}[f]$ होईल असे
    $\mModel{M}$ आणि $P(\Gamma)$ चे
    अर्थनिर्धारण~$f$ असते; विशेषतः
    $\mSat{M}{\Diamond \psi}[f(\sigma)]$.
    $\Gamma \cup
    \{\sFmla{\True}{\psi}[\sigma.n]\}$
    पूर्ततायोग्य आहे हे दाखवायचे आहे.
    यासाठी $P(\Gamma) \cup \{\sigma.n\}$
    चे अर्थनिर्धारण पुढीलप्रमाणे ठरवू:

  $\mSat{M}{\Diamond \psi}[f(\sigma)]$ असल्याने,
  $Rf(\sigma)w$ आणि $\mSat{M}{\psi}[w]$
  होईल असे $w \in W$ असते. $f'$ हे~$f$
  प्रमाणेच असू द्या, फक्त
  $f'(\sigma.n) = w$ असे ठरवा.
  $f'(\sigma) = f(\sigma)$ आणि
  $Rf(\sigma)w$ असल्याने
  $Rf'(\sigma)f'(\sigma.n)$.
  पूर्वचिन्हांच्या आरंभीच्या खंडांखालील बंदतेमुळे नव्या
  $\sigma.n$ पासून वाढवलेले कोणतेही पूर्वचिन्ह आधी वापरलेले
  नसते. त्यामुळे अर्थनिर्धारणासाठी तपासायची नवी जोडी फक्त
  $(\sigma,\sigma.n)$ हीच आहे; म्हणून $f'$
  हे $P(\Gamma) \cup \{\sigma.n\}$ चे
  अर्थनिर्धारण आहे. स्पष्टपणे
  $\mSat{M}{\psi}[f'(\sigma.n)]$.
  प्रत्येक $\sigma' \in P(\Gamma)$ साठी
  $f(\sigma') = f'(\sigma')$ असल्याने
  $\mSat{M}{\Gamma}[f']$. म्हणून
  $\mModel{M}, f'$ हे $\Gamma \cup
  \{\sFmla{\True}{\psi}[\sigma.n]\}$ पूर्ण करते.
\item $\sFmla{\False}{\Diamond \psi}[\sigma] \in \Gamma$
  याला
  $\TRule{\False}{\Diamond}$ लावून शाखा विस्तारित
  केली: शाखेवर
    $\sFmla{\False}{\psi}[\sigma.n]$ हे नवे
    चिन्हांकित सूत्र मिळते. येथे
    $\sigma.n \in P(\Gamma)$, कारण
    $\sigma.n$ वापरलेले असले पाहिजे.
    $\mSat{M}{\Gamma}[f]$ असे समजू; विशेषतः
    $\mSat/{M}{\Diamond \psi}[f(\sigma)]$.
    $f$ हे पूर्वचिन्हांचे अर्थनिर्धारण आहे आणि
    $\sigma$, $\sigma.n \in P(\Gamma)$,
    म्हणून $Rf(\sigma)f(\sigma.n)$.
    त्यामुळे $\mSat/{M}{\psi}[f(\sigma.n)]$;
    म्हणजे $\mModel{M}, f$ हे
    $\sFmla{\False}{\psi}[\sigma.n]$ पूर्ण करते.

\end{enumerate}
आता शाखा फुटणाऱ्या नियमांचा विचार करू.
\begin{enumerate}
\item $\sFmla{\False}{\psi \land \chi}[\sigma] \in \Gamma$
  याला
  $\TRule{\False}{\land}$ लावून शाखा विस्तारित
  केली. यात दोन शाखा मिळतात: डावीकडची
  $\sFmla{\False}{\psi}[\sigma]$ मधून, आणि
  उजवीकडची $\sFmla{\False}{\chi}[\sigma]$
  मधून पुढे जाते. $\mSat{M}{\Gamma}[f]$
  असे समजू; विशेषतः
  $\mSat/{M}{\psi \land \chi}[f(\sigma)]$.
  मग $\mSat/{M}{\psi}[f(\sigma)]$ किंवा
  $\mSat/{M}{\chi}[f(\sigma)]$.
  पहिल्या प्रकरणात $\mModel{M}, f$ हे
  $\sFmla{\False}{\psi}[\sigma]$ पूर्ण करते,
  म्हणून डावीकडची शाखा पूर्ततायोग्य असते.
  दुसऱ्या प्रकरणात $\mModel{M}, f$ हे
  $\sFmla{\False}{\chi}[\sigma]$ पूर्ण करते,
  म्हणून उजवीकडची शाखा पूर्ततायोग्य असते.
\item $\sFmla{\True}{\psi \lor \chi}[\sigma] \in \Gamma$
  याला
  $\TRule{\True}{\lor}$ लावून शाखा विस्तारित
  केली: सरावासाठी.
\item $\sFmla{\True}{\psi \lif \chi}[\sigma] \in \Gamma$
  याला
  $\TRule{\True}{\lif}$ लावून शाखा विस्तारित
  केली: सरावासाठी.
\end{enumerate}
\end{proof}

\begin{prob}
\ref{nml:tab:sou:thm:tableau-soundness} चा
पुरावा पूर्ण करा.
\end{prob}

\begin{cor}
\label{nml:tab:sou:cor:entailment-soundness}
$\Gamma \Proves \varphi$ असेल, तर
$\Gamma \Entails \varphi$.
\end{cor}

\begin{proof}
  $\Gamma \Proves \varphi$ असेल, तर
  $\psi_1$, \dots, $\psi_n \in \Gamma$ अशी काही सूत्रे
  असतात की $\Delta =
  \{\sFmla{\False}{\varphi}[1], \sFmla{\True}{\psi_1}[1],
  \dots, \sFmla{\True}{\psi_n}[1]\}$ साठी
  बंद टॅब्लो असतो.
  $\Gamma \Entails \varphi$ हे दाखवायचे आहे.
  तसे नाही असे समजू. मग एखाद्या
  $\mModel{M}$ आणि~$w$ साठी सर्व
  $i=1$, \dots,~$n$ बाबत
  $\mSat{M}{\psi_i}[w]$, पण
  $\mSat/{M}{\varphi}[w]$. $f(1) = w$
  असे ठेवा. मग $f$ हे~$\mModel{M}$ मधील
  $P(\Delta)$ चे अर्थनिर्धारण आहे आणि
  $\mModel{M}$ हे~$f$ च्या संदर्भात
  $\Delta$ पूर्ण करते. पण
  \ref{nml:tab:sou:thm:tableau-soundness} नुसार
  $\Delta$ साठी बंद टॅब्लो असल्यामुळे
  तो अपूर्ततायोग्य आहे. हा विरोधाभास आहे.
  म्हणून $\Gamma \Entails \varphi$.
\end{proof}

\begin{cor}
\label{nml:tab:sou:cor:weak-soundness}
$\Proves \varphi$ असेल, तर $\varphi$ सर्व
प्रतिमानांमध्ये सत्य आहे.
\end{cor}













\section{इतर प्राप्यता संबंधांसाठीचे नियम}\label{nml:tab:mru:sec}

विशिष्ट प्राप्यता संबंधांनी निर्धारित होणाऱ्या
तर्कप्रणालींसाठी \ref{nml:tab:mru:tab:more-rules} मधील
अतिरिक्त नियमांचा विचार करू.

\begin{table}
  \begin{center}
    \def\arraystretch{3}\def\fCenter{}
    \begin{tabular}{|c|c|}
    \hline

      \AxiomC{\sFmla{\True}{\Box \varphi}[\sigma]}
      \RightLabel{T$\Box$}
      \UnaryInfC{\sFmla{\True}{\varphi}[\sigma]}
      \DisplayProof
    &

      \AxiomC{\sFmla{\False}{\Diamond \varphi}[\sigma]}
      \RightLabel{T$\Diamond$}
      \UnaryInfC{\sFmla{\False}{\varphi}[\sigma]}
      \DisplayProof
    \\[1ex]
    \hline
      \AxiomC{\sFmla{\True}{\Box \varphi}[\sigma]}
      \RightLabel{D$\Box$}
      \UnaryInfC{\sFmla{\True}{\Diamond\varphi}[\sigma]}
      \DisplayProof
    &
      \AxiomC{\sFmla{\False}{\Diamond \varphi}[\sigma]}
      \RightLabel{D$\Diamond$}
      \UnaryInfC{\sFmla{\False}{\Box\varphi}[\sigma]}
      \DisplayProof
    \\[1ex]
    \hline
      \AxiomC{\sFmla{\True}{\Box \varphi}[\sigma.n]}
      \RightLabel{B$\Box$}
      \UnaryInfC{\sFmla{\True}{\varphi}[\sigma]}
      \DisplayProof
    &
      \AxiomC{\sFmla{\False}{\Diamond \varphi}[\sigma.n]}
      \RightLabel{B$\Diamond$}
      \UnaryInfC{\sFmla{\False}{\varphi}[\sigma]}
      \DisplayProof
    \\[1ex]
    \hline
      \AxiomC{\sFmla{\True}{\Box \varphi}[\sigma]}
      \RightLabel{4$\Box$}
      \UnaryInfC{\sFmla{\True}{\Box\varphi}[\sigma.n]}
      \DisplayProof
    &
      \AxiomC{\sFmla{\False}{\Diamond \varphi}[\sigma]}
      \RightLabel{4$\Diamond$}
      \UnaryInfC{\sFmla{\False}{\Diamond\varphi}[\sigma.n]}
      \DisplayProof
    \\
    $\sigma.n$ वापरलेले आहे & $\sigma.n$ वापरलेले आहे
    \\[1ex]
    \hline
      \AxiomC{\sFmla{\True}{\Box \varphi}[\sigma.n]}
      \RightLabel{4r$\Box$}
      \UnaryInfC{\sFmla{\True}{\Box\varphi}[\sigma]}
      \DisplayProof
    &
      \AxiomC{\sFmla{\False}{\Diamond \varphi}[\sigma.n]}
      \RightLabel{4r$\Diamond$}
      \UnaryInfC{\sFmla{\False}{\Diamond\varphi}[\sigma]}
      \DisplayProof
    \\[1ex]
    \hline
    \end{tabular}
  \end{center}
  \caption{अधिक मोडल नियम.}
  \label{nml:tab:mru:tab:more-rules}
\end{table}
हे नियम जोडल्यावर \ref{nml:tab:mru:tab:logics-rules} मध्ये
दिलेल्या तर्कप्रणालींसाठी निर्दोष आणि पूर्ण
अशा पद्धती मिळतात.
\begin{table}
  \begin{center}
    \begin{tabular}{lll}
      \hline
      तर्कप्रणाली & $R$ हा \dots & नियम\\
      \hline
      $\Log{T} = \Log{KT}$ & परावर्ती &
      T$\Box$
      ,
      T$\Diamond$
      \\ \hline
      $\Log{D} = \Log{KD}$ & सर्वत्र उत्तरवर्ती &
      D$\Box$
      ,
      D$\Diamond$
      \\ \hline
      $\Log{K4}$ & संक्रमक &
      4$\Box$
      ,
      4$\Diamond$
      \\ \hline
      $\Log{B} = \Log{KTB}$ & परावर्ती, &
      T$\Box$
      ,
      T$\Diamond$\\
      & सममित &
      B$\Box$
      ,
      B$\Diamond$
      \\ \hline
      $\Log{S4} = \Log{KT4}$ & परावर्ती, &
      T$\Box$
      ,
      T$\Diamond$,\\
      & संक्रमक &
      4$\Box$
      ,
      4$\Diamond$
      \\ \hline
      $\Log{S5} = \Log{KT4B}$ & परावर्ती, &
      T$\Box$
      ,
      T$\Diamond$,\\
      & संक्रमक, &
      4$\Box$
      ,
      4$\Diamond$,\\
      &  युक्लिडी &
      4r$\Box$
      ,
      4r$\Diamond$
      \\ \hline
    \end{tabular}
  \end{center}
  \caption{विविध मोडल तर्कप्रणालींसाठीचे टॅब्लो नियम.}
  \label{nml:tab:mru:tab:logics-rules}
\end{table}
\begin{ex}
  $\Log{S5} \Proves \Ax{5}$, म्हणजे
  $\Diamond\varphi \lif \Box\Diamond\varphi$, हे दाखवणारा
  बंद टॅब्लो पुढीलप्रमाणे आहे.
  \begin{oltableau}
    [\pFmla{\False}{\Diamond\formula{A} \lif \Box\Diamond \formula{A}}{1},
      just = \TAss
      [\pFmla{\True}{\Diamond \formula{A}}{1}, just = {\TRule{\False}{\lif}[1]}
        [\pFmla{\False}{\Box\Diamond \formula{A}}{1},
          just = {\TRule{\False}{\lif}[1]}
          [\pFmla{\False}{\Diamond \formula{A}}{1.1},
            just = {\TRule{\False}{\Box}[3]}
            [\pFmla{\False}{\Diamond \formula{A}}{1},
              just = {4r$\Diamond$ 4}
              [\pFmla{\True}{\formula{A}}{1.2},
                just = {\TRule{\True}{\Diamond}[2]}
                [\pFmla{\False}{\formula{A}}{1.2},
                  just = {\TRule{\False}{\Diamond}[5]}, close]
              ]
            ]
          ]
        ]
      ]
    ]
  \end{oltableau}
\end{ex}
\begin{prob}
पुढील विधाने दाखवणारे बंद टॅब्लो द्या:
  \begin{enumerate}
  \item $\Log{KT5} \Proves \Ax{B}$;
  \item $\Log{KT5} \Proves \Ax{4}$;
  \item $\Log{KDB4} \Proves \Ax{T}$;
  \item $\Log{KB4} \Proves \Ax{5}$;
  \item $\Log{KB5} \Proves \Ax{4}$;
  \item $\Log{KT} \Proves \Ax{D}$.
  \end{enumerate}
\end{prob}
\section{अतिरिक्त नियमांची निर्दोषता}\label{nml:tab:msn:sec}
प्रतिमानांच्या एखाद्या वर्गातील एका प्रतिमानात
टॅब्लोची शाखा पूर्ततायोग्य असेल, तेव्हा
नियम लावल्यावर मिळणारी शाखाही त्या वर्गातील
प्रतिमानात पूर्ततायोग्य असेल, तर तो नियम त्या
वर्गासाठी निर्दोष आहे असे म्हणतो.
\begin{prop}
  \label{nml:tab:msn:prop:soundness-T}
  \Ax{T}$\Box$
   आणि
  \Ax{T}$\Diamond$
  परावर्ती प्रतिमानांसाठी
  निर्दोष आहेत.
\end{prop}
\begin{proof}
\begin{enumerate}
\item $\sFmla{\True}{\Box \psi}[\sigma] \in \Gamma$
  याला T$\Box$ लावून शाखा विस्तारित केली:
  शाखेवर
    $\sFmla{\True}{\psi}[\sigma]$ हे नवे चिन्हांकित सूत्र मिळते. $\mSat{M}{\Gamma}[f]$
    असे समजू; विशेषतः $\mSat{M}{\Box
      \psi}[f(\sigma)]$. $R$ परावर्ती असल्याने
    $Rf(\sigma)f(\sigma)$. म्हणून
    $\mSat{M}{\psi}[f(\sigma)]$; म्हणजे
    $\mModel{M}, f$ हे
    $\sFmla{\True}{\psi}[\sigma]$ पूर्ण करते.
\item $\sFmla{\False}{\Diamond \psi}[\sigma] \in \Gamma$
  याला T$\Diamond$ लावून शाखा विस्तारित केली:
  शाखेवर
    $\sFmla{\False}{\psi}[\sigma]$ हे नवे चिन्हांकित सूत्र मिळते. $\mSat{M}{\Gamma}[f]$
    असे समजू; विशेषतः $\mSat/{M}{\Diamond
      \psi}[f(\sigma)]$. $R$ परावर्ती असल्याने
    $Rf(\sigma)f(\sigma)$. म्हणून
    $\mSat/{M}{\psi}[f(\sigma)]$; म्हणजे
    $\mModel{M}, f$ हे
    $\sFmla{\False}{\psi}[\sigma]$ पूर्ण करते.
\end{enumerate}
\end{proof}
\begin{prop}
  \label{nml:tab:msn:prop:soundness-D}
  \Ax{D}$\Box$
   आणि
  \Ax{D}$\Diamond$
  सर्वत्र उत्तरवर्ती प्रतिमानांसाठी
  निर्दोष आहेत.
\end{prop}
\begin{proof}
\begin{enumerate}
\item $\sFmla{\True}{\Box \psi}[\sigma] \in \Gamma$
  याला D$\Box$ लावून शाखा विस्तारित केली:
  शाखेवर
    $\sFmla{\True}{\Diamond \psi}[\sigma]$ हे नवे
    चिन्हांकित सूत्र मिळते. $\mSat{M}{\Gamma}[f]$
    असे समजू; विशेषतः
    $\mSat{M}{\Box \psi}[f(\sigma)]$.
    $R$ सर्वत्र उत्तरवर्ती असल्याने
    $Rf(\sigma)w$ होईल असे $w
    \in W$ असते. मग $\mSat{M}{\psi}[w]$,
    आणि म्हणून $\mSat{M}{\Diamond \psi}[f(\sigma)]$.
    त्यामुळे $\mModel{M}, f$ हे
    $\sFmla{\True}{\Diamond\psi}[\sigma]$ पूर्ण करते.
\item $\sFmla{\False}{\Diamond \psi}[\sigma] \in \Gamma$
  याला D$\Diamond$ लावून शाखा विस्तारित केली:
  शाखेवर
    $\sFmla{\False}{\Box \psi}[\sigma]$ हे नवे
    चिन्हांकित सूत्र मिळते. $\mSat{M}{\Gamma}[f]$
    असे समजू; विशेषतः
    $\mSat/{M}{\Diamond \psi}[f(\sigma)]$.
    $R$ सर्वत्र उत्तरवर्ती असल्याने
    $Rf(\sigma)w$ होईल असे $w
    \in W$ असते. मग $\mSat/{M}{\psi}[w]$,
    आणि म्हणून $\mSat/{M}{\Box \psi}[f(\sigma)]$.
    त्यामुळे $\mModel{M}, f$ हे
    $\sFmla{\False}{\Box\psi}[\sigma]$ पूर्ण करते.
\end{enumerate}
\end{proof}
\begin{prop}
  \label{nml:tab:msn:prop:soundness-B}
  \Ax{B}$\Box$
   आणि
  \Ax{B}$\Diamond$
  सममित प्रतिमानांसाठी
  निर्दोष आहेत.
\end{prop}
\begin{proof}
\begin{enumerate}
\item $\sFmla{\True}{\Box \psi}[\sigma.n] \in \Gamma$
  याला B$\Box$ लावून शाखा विस्तारित केली:
  शाखेवर
    $\sFmla{\True}{\psi}[\sigma]$ हे नवे चिन्हांकित सूत्र मिळते. $\mSat{M}{\Gamma}[f]$
    असे समजू; विशेषतः $\mSat{M}{\Box
      \psi}[f(\sigma.n)]$. $f$ हे शाखेवरील
    पूर्वचिन्हांचे~$\mModel{M}$ मधील
    अर्थनिर्धारण असल्याने
    $Rf(\sigma)f(\sigma.n)$.
    $R$ सममित असल्याने
    $Rf(\sigma.n)f(\sigma)$.
    $\mSat{M}{\Box \psi}[f(\sigma.n)]$
    असल्यामुळे $\mSat{M}{\psi}[f(\sigma)]$.
    म्हणून $\mModel{M}, f$ हे
    $\sFmla{\True}{\psi}[\sigma]$ पूर्ण करते.
\item $\sFmla{\False}{\Diamond \psi}[\sigma.n] \in \Gamma$
  याला B$\Diamond$ लावून शाखा विस्तारित केली:
  शाखेवर
    $\sFmla{\False}{\psi}[\sigma]$ हे नवे चिन्हांकित सूत्र मिळते. $\mSat{M}{\Gamma}[f]$
    असे समजू; विशेषतः $\mSat/{M}{\Diamond
      \psi}[f(\sigma.n)]$. $f$ हे शाखेवरील
    पूर्वचिन्हांचे~$\mModel{M}$ मधील
    अर्थनिर्धारण असल्याने
    $Rf(\sigma)f(\sigma.n)$.
    $R$ सममित असल्याने
    $Rf(\sigma.n)f(\sigma)$.
    $\mSat/{M}{\Diamond \psi}[f(\sigma.n)]$
    असल्यामुळे $\mSat/{M}{\psi}[f(\sigma)]$.
    म्हणून $\mModel{M}, f$ हे
    $\sFmla{\False}{\psi}[\sigma]$ पूर्ण करते.
\end{enumerate}
\end{proof}
\begin{prop}
  \label{nml:tab:msn:prop:soundness-4}
  \Ax{4}$\Box$
   आणि
  \Ax{4}$\Diamond$
  संक्रमक प्रतिमानांसाठी
  निर्दोष आहेत.
\end{prop}
\begin{proof}
\begin{enumerate}
\item $\sFmla{\True}{\Box \psi}[\sigma] \in \Gamma$
  याला 4$\Box$ लावून शाखा विस्तारित केली:
  शाखेवर
    $\sFmla{\True}{\Box\psi}[\sigma.n]$ हे नवे
    चिन्हांकित सूत्र मिळते. $\mSat{M}{\Gamma}[f]$
    असे समजू; विशेषतः
    $\mSat{M}{\Box \psi}[f(\sigma)]$.
    $f$ हे शाखेवरील पूर्वचिन्हांचे
    $\mModel{M}$ मधील अर्थनिर्धारण आहे
    आणि $\sigma.n$ वापरलेले असले पाहिजे;
    म्हणून $Rf(\sigma)f(\sigma.n)$.
    आता $Rf(\sigma.n)w$ होईल असे कोणतेही
    जग~$w$ घ्या. $R$ संक्रमक असल्याने
    $Rf(\sigma)w$. तसेच
    $\mSat{M}{\Box \psi}[f(\sigma)]$
    असल्यामुळे $\mSat{M}{\psi}[w]$.
    म्हणून $\mSat{M}{\Box \psi}[f(\sigma.n)]$
    आणि $\mModel{M}, f$ हे
    $\sFmla{\True}{\Box \psi}[\sigma.n]$ पूर्ण करते.
\item $\sFmla{\False}{\Diamond \psi}[\sigma] \in \Gamma$
  याला 4$\Diamond$ लावून शाखा विस्तारित केली:
  शाखेवर
    $\sFmla{\False}{\Diamond\psi}[\sigma.n]$ हे नवे
    चिन्हांकित सूत्र मिळते. $\mSat{M}{\Gamma}[f]$
    असे समजू; विशेषतः
    $\mSat/{M}{\Diamond \psi}[f(\sigma)]$.
    $f$ हे शाखेवरील पूर्वचिन्हांचे
    $\mModel{M}$ मधील अर्थनिर्धारण आहे
    आणि $\sigma.n$ वापरलेले असले पाहिजे;
    म्हणून $Rf(\sigma)f(\sigma.n)$.
    आता $Rf(\sigma.n)w$ होईल असे कोणतेही
    जग~$w$ घ्या. $R$ संक्रमक असल्याने
    $Rf(\sigma)w$. तसेच
    $\mSat/{M}{\Diamond \psi}[f(\sigma)]$
    असल्यामुळे $\mSat/{M}{\psi}[w]$.
    म्हणून $\mSat/{M}{\Diamond \psi}[f(\sigma.n)]$
    आणि $\mModel{M}, f$ हे
    $\sFmla{\False}{\Diamond \psi}[\sigma.n]$ पूर्ण करते.
\end{enumerate}
\end{proof}
\begin{prop}
  \label{nml:tab:msn:prop:soundness-4r}
  \Ax{4r}$\Box$
   आणि
  \Ax{4r}$\Diamond$
  युक्लिडी प्रतिमानांसाठी
  निर्दोष आहेत.
\end{prop}
\begin{proof}
\begin{enumerate}
\item $\sFmla{\True}{\Box \psi}[\sigma.n] \in \Gamma$
  याला 4r$\Box$ लावून शाखा विस्तारित केली:
  शाखेवर
    $\sFmla{\True}{\Box\psi}[\sigma]$ हे नवे
    चिन्हांकित सूत्र मिळते. $\mSat{M}{\Gamma}[f]$
    असे समजू; विशेषतः $\mSat{M}{\Box
      \psi}[f(\sigma.n)]$.
    $f$ हे शाखेवरील पूर्वचिन्हांचे
    $\mModel{M}$ मधील अर्थनिर्धारण असल्याने
    $Rf(\sigma)f(\sigma.n)$.
    आता $Rf(\sigma)w$ होईल असे कोणतेही
    जग~$w$ घ्या. $R$ युक्लिडी असल्याने
    $Rf(\sigma.n)w$.
    $\mSat{M}{\Box \psi}[f(\sigma.n)]$
    असल्याने $\mSat{M}{\psi}[w]$.
    म्हणून $\mSat{M}{\Box \psi}[f(\sigma)]$;
    आणि $\mModel{M}, f$ हे
    $\sFmla{\True}{\Box \psi}[\sigma]$ पूर्ण करते.
\item $\sFmla{\False}{\Diamond \psi}[\sigma.n] \in \Gamma$
  याला 4r$\Diamond$ लावून शाखा विस्तारित केली:
  शाखेवर
    $\sFmla{\False}{\Diamond\psi}[\sigma]$ हे नवे
    चिन्हांकित सूत्र मिळते. $\mSat{M}{\Gamma}[f]$
    असे समजू; विशेषतः $\mSat/{M}{\Diamond
      \psi}[f(\sigma.n)]$.
    $f$ हे शाखेवरील पूर्वचिन्हांचे
    $\mModel{M}$ मधील अर्थनिर्धारण असल्याने
    $Rf(\sigma)f(\sigma.n)$.
    आता $Rf(\sigma)w$ होईल असे कोणतेही
    जग~$w$ घ्या. $R$ युक्लिडी असल्याने
    $Rf(\sigma.n)w$.
    $\mSat/{M}{\Diamond \psi}[f(\sigma.n)]$
    असल्याने $\mSat/{M}{\psi}[w]$.
    म्हणून $\mSat/{M}{\Diamond \psi}[f(\sigma)]$;
    आणि $\mModel{M}, f$ हे
    $\sFmla{\False}{\Diamond \psi}[\sigma]$ पूर्ण करते.
\end{enumerate}
\end{proof}
\begin{cor}
\label{nml:tab:msn:cor:soundness-logics}
\ref{nml:tab:mru:tab:logics-rules} मध्ये दिलेल्या
टॅब्लो पद्धती त्यांच्या संबंधित
प्रतिमानवर्गांसाठी निर्दोष आहेत.
\end{cor}
\section{\Log{S5} साठी सोपे टॅब्लो}\label{nml:tab:s5:sec}
\Log{S5} ही प्रणाली वैश्विक प्रतिमानांच्या
वर्गासाठी निर्दोष आणि पूर्ण आहे; अशा प्रतिमानांत
प्रत्येक जग प्रत्येक जगापासून प्राप्य असते.
वैश्विक प्रतिमानांत प्राप्यता संबंध स्वतंत्रपणे
तपासण्याची गरज नसते: ‘$\mSat{M}{\varphi}[w]$
होईल असे जग~$w$ आहे’ हे विधान तेव्हा आणि
तेव्हाच खरे असते जेव्हा असे~$w$ हे~$u$
पासून प्राप्य असते. म्हणून \Log{S5} मध्ये
प्रतिमानाची व्याख्या केवळ जगांचा संच आणि
सत्य-मूल्यांकन~$V$ यांद्वारे करता येते. यावरून
टॅब्लोचे नियमही सोपे करता येतील असे
सुचते. सर्वसाधारण प्रकरणात पूर्वचिन्हे म्हणून
धन पूर्णांकांच्या क्रमिका वापरतो, म्हणजे कोणत्या
पूर्वचिन्हाने नाव दिलेले जग दुसऱ्या जगापासून
प्राप्य आहे याची नोंद ठेवता येते:
$\sigma.n$ हे~$\sigma$ पासून प्राप्य
जगाचे नाव देते. पण \Log{S5} मध्ये कोणतेही
जग कोणत्याही जगापासून प्राप्य असते, त्यामुळे
अशी नोंद ठेवण्याची गरज नाही. त्याऐवजी
धन पूर्णांक पूर्वचिन्हे म्हणून वापरता येतात.
सोपे केलेले नियम \ref{nml:tab:s5:tab:rules-S5}
मध्ये दिले आहेत.
\begin{table}
  \begin{center}
    \def\arraystretch{3}\def\fCenter{}
    \begin{tabular}{|c|c|}
    \hline
      \AxiomC{\sFmla{\True}{\Box \varphi}[n]}
      \RightLabel{\TRule{\True}{\Box}}
      \UnaryInfC{\sFmla{\True}{\varphi}[m]}
      \DisplayProof
      &
      \AxiomC{\sFmla{\False}{\Box \varphi}[n]}
      \RightLabel{\TRule{\False}{\Box}}
      \UnaryInfC{\sFmla{\False}{\varphi}[m]}
      \DisplayProof\\
      $m$ वापरलेले आहे & $m$ नवे आहे
      \\[1ex]
      \hline
      \AxiomC{\sFmla{\True}{\Diamond \varphi}[n]}
      \RightLabel{\TRule{\True}{\Diamond}}
      \UnaryInfC{\sFmla{\True}{\varphi}[m]}
      \DisplayProof
      &
      \AxiomC{\sFmla{\False}{\Diamond \varphi}[n]}
      \RightLabel{\TRule{\False}{\Diamond}}
      \UnaryInfC{\sFmla{\False}{\varphi}[m]}
      \DisplayProof\\
      $m$ नवे आहे & $m$ वापरलेले आहे
      \\[1ex]
      \hline
    \end{tabular}
  \end{center}
  \caption{\Log{S5} साठी सोपे केलेले नियम.}
  \label{nml:tab:s5:tab:rules-S5}
\end{table}
\begin{ex}
  $\Log{S5} \Proves \Ax{5}$, म्हणजे
  $\Diamond\varphi \lif \Box\Diamond\varphi$, हे दाखवणारा
  सोपा केलेला बंद टॅब्लो पुढीलप्रमाणे आहे.
  \begin{oltableau}
    [\pFmla{\False}{\Diamond\formula{A} \lif \Box\Diamond \formula{A}}{1},
      just = \TAss
      [\pFmla{\True}{\Diamond \formula{A}}{1}, just = {\TRule{\False}{\lif}[1]}
        [\pFmla{\False}{\Box\Diamond \formula{A}}{1},
          just = {\TRule{\False}{\lif}[1]}
          [\pFmla{\False}{\Diamond \formula{A}}{2},
            just = {\TRule{\False}{\Box}[3]}
            [\pFmla{\True}{\formula{A}}{3},
              just = {\TRule{\True}{\Diamond}[2]}
                [\pFmla{\False}{\formula{A}}{3},
                  just = {\TRule{\False}{\Diamond}[4]}, close]
            ]
          ]
        ]
      ]
    ]
  \end{oltableau}
\end{ex}
\section{\Log{K} साठी संपूर्णता}\label{nml:tab:cpl:sec}
\begin{explain}
  टॅब्लोची पद्धत संपूर्ण आहे हे
  दाखवण्यासाठी पुढील विधान सिद्ध करावे लागते:
  $\Gamma \Proves \varphi$ दाखवणारा बंद टॅब्लो
  नसेल, तर $\Gamma \Entails/ \varphi$; म्हणजे
  प्रतिदृष्टांत प्रतिमान अस्तित्वात असते.
  पण ‘बंद टॅब्लो नाही’ याचा अर्थ
  तो बांधण्याचा प्रत्येक प्रयत्न बंद होण्यात
  अयशस्वी होतो. युक्ती अशी की प्रत्येक प्रयत्न
  अयशस्वी होत असेल, तर एक विशिष्ट
  \emph{पद्धतशीर आणि सर्वसमावेशक} प्रयत्नही
  अयशस्वी होतो. बंद टॅब्लो अस्तित्वात
  असेल, तर या पद्धतशीर प्रयत्नाने तो बंद
  झाला असता. अशा एकाच टॅब्लोमधील खुल्या
  शाखांत प्रतिदृष्टांत प्रतिमानाची व्याख्या
  करण्यासाठी लागणारी सर्व माहिती असते.
  त्यातील एका खुल्या शाखेवरून मिळणाऱ्या
  प्रतिदृष्टांत प्रतिमानाची जगे म्हणजे त्या
  शाखेवर वापरलेली सर्व पूर्वचिन्हे; आणि
  विधानीय चल~$p$ हे~$\sigma$ येथे
  तेव्हा आणि तेव्हाच सत्य असते जेव्हा
  $\sFmla{\True}{p}[\sigma]$ हे त्या
  शाखेवर आलेले असते.
\end{explain}
\begin{defn}
  टॅब्लो मधील शाखेला संपूर्ण
  म्हणतात, जर नियम लावता येईल असे
  $\sFmla{S}{\varphi}[\sigma]$ हे पूर्वचिन्हयुक्त
  सूत्र शाखेवर असताना पुढील गोष्टीही
  तिच्यावर असतील:
  \begin{enumerate}
    \item विधानीय, एकाच शाखेवर सूत्रे जोडणाऱ्या
      नियमाच्या बाबतीत त्याच्या निष्कर्षांतील
      पूर्वचिन्हयुक्त सूत्रे;
    \item विधानीय, शाखा फोडणाऱ्या नियमाच्या
      बाबतीत संबंधित निष्कर्षरूप
      सूत्रांपैकी एक;
    \item नवे पूर्वचिन्ह लागणाऱ्या मोडल नियमाच्या
      बाबतीत किमान एक शक्य निष्कर्ष;
    \item वापरलेले पूर्वचिन्ह लागणाऱ्या मोडल
      नियमाच्या बाबतीत शाखेवर असलेल्या
      प्रत्येक योग्य पूर्वचिन्हासाठी संबंधित
      निष्कर्ष.
  \end{enumerate}
\end{defn}
\begin{explain}
उदाहरणार्थ, संपूर्ण शाखेवर
$\sFmla{\True}{\psi \land \chi}[\sigma]$
असेल, तर $\sFmla{\True}{\psi}[\sigma]$
आणि $\sFmla{\True}{\chi}[\sigma]$
ही दोन्ही असतात. तिच्यावर
$\sFmla{\True}{\psi \lor \chi}[\sigma]$
असेल, तर $\sFmla{\True}{\psi}[\sigma]$
आणि $\sFmla{\True}{\chi}[\sigma]$
यांपैकी किमान एक असते.
शाखेवर $\sFmla{\False}{\Box \psi}[\sigma]$
असेल, तर काही~$n$ साठी अनुक्रमे
$\sFmla{\False}{\psi}[\sigma.n]$
हेही असते. आणि शाखेवर
$\sFmla{\True}{\Box \psi}[\sigma]$
असेल, तर शाखेवर $\sigma.n$ वापरलेले
असेल अशा प्रत्येक~$n$ साठी अनुक्रमे
$\sFmla{\True}{\psi}[\sigma.n]$
हेही असते.
\end{explain}
\begin{prop}\label{nml:tab:cpl:prop:complete-tableau}
  $\Gamma$ हा $1$ या एकाच पूर्वचिन्हाने चिन्हांकित सूत्रांचा
  सांत संच असू द्या. त्यासाठी असा टॅब्लो असतो की
  त्याची प्रत्येक शाखा बंद किंवा संपूर्ण असते.
\end{prop}
\begin{proof}
  $\Gamma$ साठीच्या टॅब्लो मधील
  एखादी खुली शाखा घ्या. त्या शाखेवर नियम
  लावता येतील अशी सांत पूर्वचिन्हयुक्त
  सूत्रे आहेत. ठरलेल्या एखाद्या क्रमाने
  (उदाहरणार्थ, वरून खाली), ज्या
  पूर्वचिन्हयुक्त सूत्रे साठी
  (1)--(4) या अटी आधीच पूर्ण होत नाहीत,
  त्यांपैकी प्रत्येकाला लागू होणारे नियम
  लावून शाखा विस्तारित करा. काही वेळा शाखा
  फुटतील; मग उरलेल्या पूर्वचिन्हयुक्त
  सूत्रे साठी प्रत्येक नव्या शाखेच्या
  टोकावर नियम लावा. पूर्वचिन्हयुक्त
  सूत्रे आणि शाखेवर वापरलेली
  पूर्वचिन्हे यांची संख्या सांत असल्याने
  या प्रक्रियेतून मूळ शाखेपासून पुढे जाणाऱ्या
  (कदाचित अनेक) शाखा मिळतात. त्यांपैकी
  प्रत्येकावर ही प्रक्रिया पुन्हा लावा.
  बंद झालेल्या शाखांवर पुढील नियम लावू नयेत.
  ही प्रक्रिया सांत पायऱ्यांत थांबते: नव्या पूर्वचिन्हावर येणारी
  सूत्रे त्याच्या पालकावरील मोडल सूत्रांची उपसूत्रे असतात.
  त्यामुळे प्रत्येक पूर्वचिन्ह-विस्ताराने उपलब्ध मोडल खोली कमी
  होते; प्रारंभीच्या सांत सूत्रसंचाची कमाल मोडल खोली ही शाखेवरील
  पूर्वचिन्हांच्या उंचीची मर्यादा ठरते. प्रत्येक पूर्वचिन्हावर फक्त
  प्रारंभीच्या सूत्रांची सांत चिन्हांकित उपसूत्रे येऊ शकतात.
  प्रत्येक नवेपणा मागणाऱ्या चिन्हांकित मोडल सूत्रासाठी एकच
  साक्षीदार पूर्वचिन्ह पुरते, म्हणून प्रत्येक पूर्वचिन्हाची संततीही
  सांत असते. मर्यादित उंची आणि प्रत्येक ठिकाणी सांत संतती
  यांमुळे शाखेवरील सर्व पूर्वचिन्हे सांत संख्येची असतात.
  उर्वरित प्रत्येक नियमाचे आवश्यक अनुप्रयोगही सांत असतात.
  रचनेमुळे उरलेली प्रत्येक खुली शाखा संपूर्ण असते.
\end{proof}
\begin{thm}[संपूर्णता]
  \label{nml:tab:cpl:thm:tableau-completeness}
  $\Gamma$ हा $1$ या एकाच पूर्वचिन्हाने चिन्हांकित सूत्रांचा संच असू द्या.
  $\Gamma$ साठी बंद टॅब्लो नसेल,
  तर $\Gamma$ पूर्ततायोग्य आहे.
\end{thm}
\begin{proof}
$\Gamma$ रिक्त असेल, तर एक जग असलेले कोणतेही प्रतिमान
ते पूर्ण करते. प्रथम $\Gamma$ रिक्त नसलेला सांत संच आहे
असे समजू. वरील विधानानुसार, $\Gamma$ साठी असा
टॅब्लो असतो की त्याची प्रत्येक शाखा
बंद किंवा संपूर्ण असते. बंद टॅब्लो नसल्याने,
त्या टॅब्लोमध्ये किमान एक शाखा
खुली आणि संपूर्ण असते. त्या शाखेवरील
पूर्वचिन्हयुक्त सूत्रांचा संच
$\Delta$ आणि त्यात आलेल्या पूर्वचिन्हांचा
संच $P(\Delta)$ असू द्या.
आता $\mModel{M(\Delta)} = \tuple{P(\Delta), R, V}$
हे प्रतिमान ठरवू. त्यातील जगे म्हणजे~$\Delta$
मधील पूर्वचिन्हे, आणि प्राप्यता संबंध असा आहे:
\[
R\sigma\sigma' \quad \text{तेव्हा आणि तेव्हाच} \quad
\sigma'=\sigma.n \quad \text{काही~$n$ साठी}
\]
आणि
\[V(p) = \Setabs{\sigma}{\sFmla{\True}{p}[\sigma] \in \Delta}.
\]
आता~$\varphi$ च्या रचनेवरून आगमनाने पुढील
दाखवू: $\sFmla{\True}{\varphi}[\sigma] \in
\Delta$ असेल, तर
$\mSat{M(\Delta)}{\varphi}[\sigma]$; आणि
$\sFmla{\False}{\varphi}[\sigma] \in \Delta$
असेल, तर $\mSat/{M(\Delta)}{\varphi}[\sigma]$.

\begin{enumerate}
  \item \indcase{\varphi}{p}{$\sFmla{\True}{\indfrm}[\sigma] \in \Delta$
    असेल, तर~$V$ च्या व्याख्येनुसार
    $\sigma \in V(p)$; म्हणून
    $\mSat{M(\Delta)}{\indfrm}[\sigma]$.

    $\sFmla{\False}{\indfrm}[\sigma] \in \Delta$ असेल, तर
    $\sFmla{\True}{\indfrm}[\sigma] \notin \Delta$,
    कारण अन्यथा शाखा बंद झाली असती.
    म्हणून $\sigma \notin V(p)$, आणि
    $\mSat/{M(\Delta)}{\indfrm}[\sigma]$.}
  \item \indcase{\varphi}{\lnot \psi}{जर
      $\sFmla{\True}{\indfrm}[\sigma] \in \Delta$,
      तर शाखा संपूर्ण असल्याने
      $\sFmla{\False}{\psi}[\sigma] \in \Delta$.
      आगमनाच्या गृहीतकानुसार
      $\mSat/{M(\Delta)}{\psi}[\sigma]$;
      म्हणून $\mSat{M(\Delta)}{\indfrm}[\sigma]$.

      $\sFmla{\False}{\indfrm}[\sigma] \in \Delta$
      असेल, तर शाखा संपूर्ण असल्याने
      $\sFmla{\True}{\psi}[\sigma] \in \Delta$.
      आगमनाच्या गृहीतकानुसार
      $\mSat{M(\Delta)}{\psi}[\sigma]$;
      म्हणून $\mSat/{M(\Delta)}{\indfrm}[\sigma]$.}
  \item \indcase{\varphi}{\psi \land \chi}{जर
      $\sFmla{\True}{\indfrm}[\sigma] \in \Delta$,
      तर शाखा संपूर्ण असल्याने
      $\sFmla{\True}{\psi}[\sigma] \in \Delta$
      आणि $\sFmla{\True}{\chi}[\sigma] \in \Delta$
      ही दोन्ही. आगमनाच्या गृहीतकानुसार
      $\mSat{M(\Delta)}{\psi}[\sigma]$ आणि
      $\mSat{M(\Delta)}{\chi}[\sigma]$.
      म्हणून $\mSat{M(\Delta)}{\indfrm}[\sigma]$.

      $\sFmla{\False}{\indfrm}[\sigma] \in \Delta$
      असेल, तर शाखा संपूर्ण असल्याने
      $\sFmla{\False}{\psi}[\sigma] \in \Delta$
      किंवा $\sFmla{\False}{\chi}[\sigma] \in \Delta$.
      आगमनाच्या गृहीतकानुसार
      $\mSat/{M(\Delta)}{\psi}[\sigma]$ किंवा
      $\mSat/{M(\Delta)}{\chi}[\sigma]$.
      म्हणून $\mSat/{M(\Delta)}{\indfrm}[\sigma]$.}
  \item \indcase{\varphi}{\psi \lor \chi}{जर
      $\sFmla{\True}{\indfrm}[\sigma] \in \Delta$,
      तर शाखा संपूर्ण असल्याने
      $\sFmla{\True}{\psi}[\sigma] \in \Delta$
      किंवा $\sFmla{\True}{\chi}[\sigma] \in \Delta$.
      आगमनाच्या गृहीतकानुसार
      $\mSat{M(\Delta)}{\psi}[\sigma]$ किंवा
      $\mSat{M(\Delta)}{\chi}[\sigma]$.
      म्हणून $\mSat{M(\Delta)}{\indfrm}[\sigma]$.

      $\sFmla{\False}{\indfrm}[\sigma] \in \Delta$
      असेल, तर शाखा संपूर्ण असल्याने
      $\sFmla{\False}{\psi}[\sigma] \in \Delta$
      आणि $\sFmla{\False}{\chi}[\sigma] \in \Delta$
      ही दोन्ही. आगमनाच्या गृहीतकानुसार
      $\mSat/{M(\Delta)}{\psi}[\sigma]$ आणि
      $\mSat/{M(\Delta)}{\chi}[\sigma]$.
      म्हणून $\mSat/{M(\Delta)}{\indfrm}[\sigma]$.}
  \item \indcase{\varphi}{\psi \lif \chi}{जर
      $\sFmla{\True}{\indfrm}[\sigma] \in \Delta$,
      तर शाखा संपूर्ण असल्याने
      $\sFmla{\False}{\psi}[\sigma] \in \Delta$
      किंवा $\sFmla{\True}{\chi}[\sigma] \in \Delta$.
      आगमनाच्या गृहीतकानुसार
      $\mSat/{M(\Delta)}{\psi}[\sigma]$ किंवा
      $\mSat{M(\Delta)}{\chi}[\sigma]$.
      म्हणून $\mSat{M(\Delta)}{\indfrm}[\sigma]$.

      $\sFmla{\False}{\indfrm}[\sigma] \in \Delta$
      असेल, तर शाखा संपूर्ण असल्याने
      $\sFmla{\True}{\psi}[\sigma] \in \Delta$
      आणि $\sFmla{\False}{\chi}[\sigma] \in \Delta$
      ही दोन्ही. आगमनाच्या गृहीतकानुसार
      $\mSat{M(\Delta)}{\psi}[\sigma]$ आणि
      $\mSat/{M(\Delta)}{\chi}[\sigma]$.
      म्हणून $\mSat/{M(\Delta)}{\indfrm}[\sigma]$.}
  \item \indcase{\varphi}{\Box \psi}{जर
      $\sFmla{\True}{\indfrm}[\sigma] \in \Delta$,
      तर शाखा संपूर्ण असल्याने तिच्यावर
      वापरलेल्या प्रत्येक $\sigma.n$ साठी
      $\sFmla{\True}{\psi}[\sigma.n] \in \Delta$.
      म्हणजे $R\sigma\sigma'$ होईल अशा
      प्रत्येक $\sigma' \in P(\Delta)$ साठी
      हे खरे आहे. आगमनाच्या गृहीतकानुसार,
      $R\sigma\sigma'$ असलेल्या प्रत्येक
      $\sigma'$ साठी
      $\mSat{M(\Delta)}{\psi}[\sigma']$.
      म्हणून $\mSat{M(\Delta)}{\indfrm}[\sigma]$.

      $\sFmla{\False}{\indfrm}[\sigma] \in \Delta$
      असेल, तर शाखा संपूर्ण असल्याने
      काही $\sigma.n$ साठी
      $\sFmla{\False}{\psi}[\sigma.n] \in \Delta$.
      आगमनाच्या गृहीतकानुसार
      $\mSat/{M(\Delta)}{\psi}[\sigma.n]$.
      $R\sigma(\sigma.n)$ असल्याने
      $\mSat/{M(\Delta)}{\psi}[\sigma']$
      होईल असे~$\sigma'$ असते. म्हणून
      $\mSat/{M(\Delta)}{\indfrm}[\sigma]$.}
  \item \indcase{\varphi}{\Diamond \psi}{जर
      $\sFmla{\True}{\indfrm}[\sigma] \in \Delta$,
      तर शाखा संपूर्ण असल्याने
      काही $\sigma.n$ साठी
      $\sFmla{\True}{\psi}[\sigma.n] \in \Delta$.
      आगमनाच्या गृहीतकानुसार
      $\mSat{M(\Delta)}{\psi}[\sigma.n]$.
      $R\sigma(\sigma.n)$ असल्याने
      $\mSat{M(\Delta)}{\psi}[\sigma']$
      होईल असे~$\sigma'$ असते. म्हणून
      $\mSat{M(\Delta)}{\indfrm}[\sigma]$.

      $\sFmla{\False}{\indfrm}[\sigma] \in \Delta$
      असेल, तर शाखा संपूर्ण असल्याने तिच्यावर
      वापरलेल्या प्रत्येक $\sigma.n$ साठी
      $\sFmla{\False}{\psi}[\sigma.n] \in \Delta$.
      म्हणजे $R\sigma\sigma'$ होईल अशा
      प्रत्येक $\sigma' \in P(\Delta)$ साठी
      हे खरे आहे. आगमनाच्या गृहीतकानुसार,
      $R\sigma\sigma'$ असलेल्या प्रत्येक
      $\sigma'$ साठी
      $\mSat/{M(\Delta)}{\psi}[\sigma']$.
      म्हणून $\mSat/{M(\Delta)}{\indfrm}[\sigma]$.}
\end{enumerate}
$\Gamma \subseteq \Delta$ असल्याने, प्रत्येक पूर्वचिन्हाला तेच जग
नेमणारे अर्थनिर्धारण घेऊन $\mSat{M(\Delta)}{\Gamma}[\mathrm{id}]$.
याने सांत प्रकरण सिद्ध झाले.
असांत $\Gamma$ साठी त्याच्या प्रत्येक सांत उपसंचाला बंद टॅब्लो
नसतो: तसा टॅब्लो असता, तर त्याची सांत गृहीतके $\Gamma$
मध्येही असल्याने तो $\Gamma$ साठी बंद टॅब्लो ठरला असता.
म्हणून सांत प्रकरणावरून प्रत्येक सांत उपसंच पूर्ततायोग्य आहे.
चिन्हांकित गृहीतकांचा साध्या मोडल सूत्रांचा संच असा ठरवू:
\[\Upsilon = \Setabs{\psi}{\sFmla{\True}{\psi}[1]\in\Gamma}
\cup \Setabs{\lnot\psi}{\sFmla{\False}{\psi}[1]\in\Gamma}.
\]
याच्या प्रत्येक सांत उपसंचाला मूळ $\Gamma$ चा सांत उपसंच
पूर्ण करणाऱ्या प्रतिमानात एकाच जगावर सत्य करता येते.
त्यामुळे $\Upsilon$ हा \Log{K}-सुसंगत संच आहे: त्यावरून
असत्य सूत्र निष्पन्न होत असते, तर त्या सांत सिद्धतेतील सांत
गृहीतकांच्या पूर्ततायोग्यतेला \Log{K} ची निर्दोषता विरोधी ठरली
असती. \ref{nml:com:lin:thm:lindenbaum} नुसार
$\Upsilon$ ला संपूर्ण \Log{K}-सुसंगत संच $\Lambda$ पर्यंत
वाढवता येते. \ref{nml:com:tru:prop:truthlemma} नुसार
कॅनॉनिकल प्रतिमानात $\Lambda$ या जगावर $\Upsilon$ मधील
प्रत्येक सूत्र सत्य असते. $f(1)=\Lambda$ हे अर्थनिर्धारण मूळ
चिन्हांकित संच $\Gamma$ पूर्ण करते. याने असांत प्रकरणही सिद्ध झाले.
\end{proof}

\begin{prob}
\ref{nml:tab:cpl:thm:tableau-completeness} चा
पुरावा पूर्ण करा.
\end{prob}

\begin{cor}
\label{nml:tab:cpl:cor:entailment-completeness}
$\Gamma \Entails \varphi$ असेल, तर
$\Gamma \Proves \varphi$.
\end{cor}

\begin{cor}
\label{nml:tab:cpl:cor:weak-completeness}
$\varphi$ सर्व प्रतिमानांमध्ये सत्य असेल,
तर $\Proves \varphi$.
\end{cor}













\section{टॅब्लो मधून प्रतिदृष्टांत प्रतिमाने}\label{nml:tab:cou:sec}

\begin{explain}
  संपूर्णता प्रमेयाचा पुरावा
  $\Entails \varphi$ असेल तर $\Proves \varphi$
  हे दाखवतोच; शिवाय $\Entails/ \varphi$
  असेल, तर~$\varphi$ साठी प्रतिदृष्टांत
  प्रतिमान बांधण्याची पद्धतही देतो.
  $\Log{K}$ साठी ही पद्धत
  \emph{निर्णय प्रक्रिया} आहे. कारण
  $\Entails/ \varphi$ असे समजू. मग
  \ref{nml:tab:cpl:prop:complete-tableau} चा
  पुरावा संपूर्ण टॅब्लो बांधण्याची
  पद्धत देतो. ती नेहमी थांबतेही.
  $\Log{K}$ चे विधानीय नियम कमी गुंतागुंतीची
  पूर्वचिन्हयुक्त सूत्रे जोडतात;
  म्हणजे प्रत्येक
  $\sFmla{S}{\varphi}[\sigma]$ या चिन्हांकित
  सूत्रासाठी प्रत्येक विधानीय नियम
  शाखेवर एकदाच लावावा लागतो.
  नवी पूर्वचिन्हे फक्त
  $\TRule{\False}{\Box}$
    आणि $\TRule{\True}{\Diamond}$
  या नियमांनी
  तयार होतात. शाखेवरील प्रत्येक पूर्वचिन्हयुक्त
  चिन्हांकित मोडल सूत्रासाठी संबंधित नियम
  एकदाच लावावा लागतो; त्या अनुप्रयोगात
  एकच नवे पूर्वचिन्ह मिळते.
  $\TRule{\True}{\Box}$
    आणि $\TRule{\False}{\Diamond}$
  हे नियम
  अनेकदा लावावे लागू शकतात; पण प्रत्येक
  पूर्वचिन्हयुक्त चिन्हांकित मोडल सूत्रासाठी
  त्याच्या पूर्वचिन्हाचा शाखेवर वापरलेला प्रत्येक
  तात्काळ विस्तार घेऊन एकदाच लावणे पुरते.
  \ref{nml:tab:cpl:prop:complete-tableau} च्या पुराव्यात
  दाखवल्याप्रमाणे, सांत प्रारंभिक सूत्रसंचाची मोडल
  खोली पूर्वचिन्हांची उंची मर्यादित करते; प्रत्येक
  पूर्वचिन्हाची संततीही सांत असते. म्हणून नवी
  पूर्वचिन्हे सांत संख्येनेच निर्माण होतात. म्हणून
  सांत पायऱ्यांनंतर शाखा बंद होते किंवा
  संपूर्ण होते.

  खुली संपूर्ण शाखा असलेला टॅब्लो
  बांधल्यावर \ref{nml:tab:cpl:thm:tableau-completeness}
  चा पुरावा मूळ पूर्वचिन्हयुक्त सूत्रे
  चा संच पूर्ण करणारे स्पष्ट प्रतिमान देतो.
  म्हणून $\Gamma \Entails \varphi$ असेल, तर
  बंद टॅब्लो अस्तित्वात असतो आणि
  $\Gamma \Proves \varphi$. याशिवाय, $\Gamma$ सांत
  असेल, तर योग्य पद्धतीने बंद टॅब्लो शोधताना
  त्याऐवजी ‘संपूर्ण’ टॅब्लो मिळाला,
  तर $\Gamma \Entails/ \varphi$ हे कळते
  आणि प्रतिदृष्टांत प्रतिमान प्रत्यक्ष
  बांधताही येते.
\end{explain}


\begin{ex}
  $\Proves/ \Box(p \lor q) \lif (\Box p \lor \Box
  q)$ हे आपल्याला माहीत आहे. टॅब्लो बांधण्याची
  सुरुवात पुढीलप्रमाणे होते:
  \begin{oltableau}
    [\pFmla{\False}{\Box(p \lor q) \lif (\Box p \lor \Box q)}{1},
      just = \TAss, checked
      [\pFmla{\True}{\Box(p \lor q)}{1},
        just = {\TRule{\False}{\lif}[1]},
        [\pFmla{\False}{\Box p \lor \Box q}{1},
          just = {\TRule{\False}{\lif}[1]}, checked
          [\pFmla{\False}{\Box p}{1},
            just = {\TRule{\False}{\lor}[3]}, checked
            [\pFmla{\False}{\Box q}{1},
              just = {\TRule{\False}{\lor}[3]}, checked
              [\pFmla{\False}{p}{1.1},
                just = {\TRule{\False}{\Box}[4]}, checked
                [\pFmla{\False}{q}{1.2},
                  just = {\TRule{\False}{\Box}[5]}, checked
                ]
              ]
            ]
          ]
        ]
      ]
    ]
  \end{oltableau}
  अर्थात हा टॅब्लो अजून पूर्ण झालेला नाही.
  पुढे खूण नसलेल्या एकमेव ओळीचा विचार करू:
  ओळ~$2$ वरील पूर्वचिन्हयुक्त सूत्र
  $\sFmla{\True}{\Box(p \lor q)}[1]$.
  पद्धतशीर टॅब्लो बांधताना शाखेवर वापरलेल्या
  $1.n$ या आकाराच्या प्रत्येक पूर्वचिन्हासाठी
  $\TRule{\True}{\Box}$ लावायचा असतो;
  म्हणजे $1.1$ आणि $1.2$ या दोन्हींसाठी:
  \begin{oltableau}
    [\pFmla{\False}{\Box(p \lor q) \lif (\Box p \lor \Box q)}{1},
      just = \TAss, checked
      [\pFmla{\True}{\Box(p \lor q)}{1},
        just = {\TRule{\False}{\lif}[1]},
        [\pFmla{\False}{\Box p \lor \Box q}{1},
          just = {\TRule{\False}{\lif}[1]}, checked
          [\pFmla{\False}{\Box p}{1},
            just = {\TRule{\False}{\lor}[3]}, checked
            [\pFmla{\False}{\Box q}{1},
              just = {\TRule{\False}{\lor}[3]}, checked
              [\pFmla{\False}{p}{1.1},
                just = {\TRule{\False}{\Box}[4]}, checked
                [\pFmla{\False}{q}{1.2},
                  just = {\TRule{\False}{\Box}[5]}, checked
                  [\pFmla{\True}{p \lor q}{1.1},
                    just = {\TRule{\True}{\Box}[2]}
                    [\pFmla{\True}{p \lor q}{1.2},
                      just = {\TRule{\True}{\Box}[2]}
                    ]
                  ]
                ]
              ]
            ]
          ]
        ]
      ]
    ]
  \end{oltableau}
  आता ओळी 2, 8 आणि 9 यांवर खूण नाही.
  पण नवे पूर्वचिन्ह जोडलेले नाही. म्हणून
  ओळी~8 आणि~9 वर
  $\TRule{\True}{\lor}$ लावू; तयार होणाऱ्या
  सर्व शाखांवर, त्या बंद होत नसतील तोवर,
  हाच नियम लागू करू:
  \begin{oltableau}
    [\pFmla{\False}{\Box(p \lor q) \lif (\Box p \lor \Box q)}{1},
      just = \TAss, checked
      [\pFmla{\True}{\Box(p \lor q)}{1},
        just = {\TRule{\False}{\lif}[1]}, checked
        [\pFmla{\False}{\Box p \lor \Box q}{1},
          just = {\TRule{\False}{\lif}[1]}, checked
          [\pFmla{\False}{\Box p}{1},
            just = {\TRule{\False}{\lor}[3]}, checked
            [\pFmla{\False}{\Box q}{1},
              just = {\TRule{\False}{\lor}[3]}, checked
              [\pFmla{\False}{p}{1.1},
                just = {\TRule{\False}{\Box}[4]}, checked
                [\pFmla{\False}{q}{1.2},
                  just = {\TRule{\False}{\Box}[5]}, checked
                  [\pFmla{\True}{p \lor q}{1.1},
                    just = {\TRule{\True}{\Box}[2]}, checked
                    [\pFmla{\True}{p \lor q}{1.2},
                      just = {\TRule{\True}{\Box}[2]}, checked
                      [\pFmla{\True}{p}{1.1},
                        just = {\TRule{\True}{\lor}[8]}, checked, close
                      ]
                      [\pFmla{\True}{q}{1.1},
                        just = {\TRule{\True}{\lor}[8]}, checked
                        [\pFmla{\True}{p}{1.2},
                          just = {\TRule{\True}{\lor}[9]}, checked]
                        [\pFmla{\True}{q}{1.2},
                          just = {\TRule{\True}{\lor}[9]}, checked, close]
                      ]
                    ]
                  ]
                ]
              ]
            ]
          ]
        ]
      ]
    ]
  \end{oltableau}
  आता एकच खुली शाखा उरते; ती संपूर्ण आहे.
  तिच्यावरून प्रतिमान असे ठरवू: जगे
  $W = \{1, 1.1, 1.2\}$ (खुल्या शाखेवर
  येणारी हीच पूर्वचिन्हे), प्राप्यता संबंध
  $R = \{\tuple{1, 1.1}, \tuple{1, 1.2}\}$,
  आणि सत्य-मूल्यांकन $V(p) = \{1.2\}$
  (कारण ओळ~11 वर
  $\sFmla{\True}{p}[1.2]$ आहे) व
  $V(q) = \{1.1\}$ (कारण ओळ~10 वर
  $\sFmla{\True}{q}[1.1]$ आहे).
  \ref{nml:tab:cou:fig:counter-Box} मध्ये हे प्रतिमान
  दाखवले आहे. ते $\Box(p \lor q) \lif
  (\Box p \lor \Box q)$ चे प्रतिदृष्टांत
  प्रतिमान आहे हे तपासता येईल.
  \begin{figure}
  \begin{center}
    \begin{tikzpicture}[modal]
      \node[world] (w1) [label={[align=right]right:\mFalse{p}\\ \mFalse{q}}]{$1$};
      \node[world] (w2) [label={[align=right]right:\mFalse{p}\\ \mTrue{q}},
        above left=of w1]{$1.1$};
      \node[world] (w3) [label={[align=right]right:\mTrue{p}\\ \mFalse{q}},
        above right=of w1] {$1.2$};
      \draw[->] (w1) to (w2);
      \draw[->] (w1) to (w3);
    \end{tikzpicture}
  \end{center}
  \caption{$\Box(p \lor q) \lif (\Box p \lor \Box
    q)$ चे प्रतिदृष्टांत प्रतिमान.}
  \label{nml:tab:cou:fig:counter-Box}
\end{figure}
\end{ex}




